% !TEX root = Tesi_Corsi_Leonardo_656276.tex
% Preambolo condiviso fra la tesi completa e il PDF dei soli capitoli 1-3.
% Non contiene \begin{document}: lo apre corpo.tex.

% Template Tesi in Informatica, ispirata al template "Tesi di laurea - Università di Pisa" di Simone Schirinzi (https://www.overleaf.com/latex/templates/tesi-di-laurea-universita-di-pisa/rwdcqtqwftpg).

% Carattere dimensione 12
\documentclass[12pt]{report}

% Per la stampa fronte-retro sostituire con:
% \documentclass[12pt, twoside]{report}

% Margini (4cm a sx, 2.5cm a dx, 2.5cm in alto, 2.5cm in basso)
% head=21.75pt: spazio per la testatina, il valore richiesto da scrlayer-scrpage
\usepackage[top=2.5cm, bottom=2.5cm, left=4cm, right=2.5cm, head=21.75pt, centering]{geometry}

% Per la stampa fronte-retro sostituire con: 
% \usepackage[top=2.5cm, bottom=2.5cm, inner=4cm, outer=4cm, right=2.5cm, centering]{geometry}

% Interlinea
\linespread{1.5}

% Librerie utili
\usepackage[italian]{babel} % applicazione regole di scrittura per la lingua italiana 
\usepackage[utf8]{inputenc} % codifica UTF-8
\usepackage{scrlayer-scrpage} % stili pagina: testatine e pie' di pagina
% Testatina: capitolo a sinistra, numero di pagina a destra, riga di separazione.
% Le pagine di apertura capitolo restano senza testatina (stile plain) e il
% numero di pagina lo tengono in fondo: e' la forma fra parentesi quadre.
\automark[section]{chapter}
\ihead[]{\headmark}
\chead[]{}
\ohead[]{\pagemark}
\ifoot[]{}
\cfoot[]{}
\ofoot[\pagemark]{}
\KOMAoptions{headsepline=true}
\pagestyle{scrheadings}
\usepackage{mathptmx} % font Times New Roman (simile)
\usepackage{graphicx} % inserimento di immagini
\usepackage{csquotes} % per le citazioni "in blocco"
\usepackage{url} % per \url{} nelle voci di bibliografia
%\usepackage[backend=biber, sorting=nty, ]{biblatex} % bibliografia con pacchetto biblatex (https://ctan.org/pkg/biblatex?lang=en)
\appto{\bibsetup}{\raggedright}

\usepackage{titlesec} % per la formattazione dei titoli delle sezioni, capitoli etc.
\usepackage{float} % per il posizionamento delle immagini
\usepackage{booktabs} % \toprule \midrule \bottomrule per le tabelle dei risultati
\usepackage{enumitem} % spaziatura degli elenchi: \begin{itemize}[topsep=..., itemsep=...]
\usepackage{needspace} % \needspace{N\baselineskip}: non spezzare un blocco a fine pagina
\usepackage{tikz} % per i disegni dell'architettura
\usetikzlibrary{arrows.meta, calc}
\usepackage{pgfplots} % per i grafici delle misure
\pgfplotsset{compat=1.18}
\usepackage{amsmath} % per le formule dei costi
\usepackage{amssymb} % \square a fine dimostrazione

\usepackage{listings} % per il codice di programmazione
% Fonte https://en.wikibooks.org/wiki/LaTeX/Source_Code_Listings. Per la lista di sintassi riconosciute.
\renewcommand{\lstlistingname}{Codice}% Listing -> Codice
\usepackage{xcolor}  % stile del codice
\definecolor{mygreen}{rgb}{0,0.6,0}
\definecolor{mygray}{rgb}{0.5,0.5,0.5}
\definecolor{mymauve}{rgb}{0.58,0,0.82}
\definecolor{darkgray}{rgb}{.4,.4,.4}
\definecolor{navy}{HTML}{000080}
\definecolor{purple}{rgb}{0.65, 0.12, 0.82}
\definecolor{codepurple}{rgb}{0.58,0,0.82}
\definecolor{backcolour}{rgb}{0.95,0.95,0.92}

% Stili configurabili del codice (lslisting) 
\lstset{ %
belowcaptionskip=0.5em,
backgroundcolor=\color{backcolour}, % choose the background color; you must add \usepackage{color} or \usepackage{xcolor}
basicstyle=\footnotesize, % the size of the fonts that are used for the code
breakatwhitespace=false, % sets if automatic breaks should only happen at whitespace
breaklines=true, % sets automatic line breaking
captionpos=b, % sets the caption-position to bottom
commentstyle=\color{mygreen}, % comment style
deletekeywords={...}, % if you want to delete keywords from the given language
escapeinside={\%*}{*)}, % if you want to add LaTeX within your code
extendedchars=true, % lets you use non-ASCII characters; for 8-bits encodings only, does not work with UTF-8
frame=single, % adds a frame around the code
keepspaces=true, % keeps spaces in text, useful for keeping indentation of code (possibly needs columns=flexible)
keywordstyle=\color{codepurple}, % keyword style
% language=Octave, % the language of the code
morekeywords={*,...}, % if you want to add more keywords to the set
numbers=left, % where to put the line-numbers; possible values are (none, left, right)
numbersep=5pt, % how far the line-numbers are from the code
numberstyle=\tiny\color{mygray}, % the style that is used for the line-numbers
rulecolor=\color{black}, % if not set, the frame-color may be changed on line-breaks within not-black text (e.g. comments (green here))
showspaces=false, % show spaces everywhere adding particular underscores; it overrides 'showstringspaces'
showstringspaces=false, % underline spaces within strings only
showtabs=false, % show tabs within strings adding particular underscores
stepnumber=1, % the step between two line-numbers. If it's 1, each line will be numbered
stringstyle=\color{mymauve}, % string literal style
tabsize=2, % sets default tabsize to 2 spaces
title=\lstname % show the filename of files included with \lstinputlisting; also try caption instead of title
}

% END of listing package 

% Esempio riconoscimento sintassi JavaScript 

\lstdefinelanguage{JavaScript}{
  keywords={typeof, new, true, false, catch, function, return, null, catch, switch, var, const, let, if, in, while, do, else, case, break},
  keywordstyle=\color{blue}\bfseries,
  ndkeywords={class, export, boolean, throw, implements, import, this},
  ndkeywordstyle=\color{darkgray}\bfseries,
  identifierstyle=\color{black},
  sensitive=false,
  comment=[l]{//},
  morecomment=[s]{/*}{*/},
  commentstyle=\color{mygreen}\ttfamily,
  stringstyle=\color{red}\ttfamily,
  morestring=[b]',
  morestring=[b]"
}

% Assembly RISC-V: le istruzioni rv32i usate dai kernel, piu' le quattro
% istruzioni di comunicazione dell'estensione, evidenziate a parte.
\lstdefinelanguage{RISCV}{
  keywords={add, addi, sub, and, andi, or, ori, xor, xori, sll, slli, srl, srli,
            sra, srai, slt, slti, sltu, sltiu, lui, auipc, lw, lh, lhu, lb, lbu,
            sw, sh, sb, beq, bne, blt, bge, bltu, bgeu, beqz, bnez, blez, bgez,
            j, jal, jalr, mv, li, la, ret, nop, ecall, mul},
  keywordstyle=\color{codepurple}\bfseries,
  ndkeywords={IN, OUT, ISRDY, SETRDY, NORD, EST, SUD, OVEST},
  ndkeywordstyle=\color{navy}\bfseries,
  identifierstyle=\color{black},
  sensitive=true,
  comment=[l]{\#},
  commentstyle=\color{mygreen}\ttfamily,
  morestring=[b]",
  stringstyle=\color{mymauve}\ttfamily
}

\lstset{
   extendedchars=true,
   basicstyle=\footnotesize\ttfamily,
   showstringspaces=false,
   showspaces=false,
   numbers=left,
   numberstyle=\footnotesize,
   numbersep=9pt,
   tabsize=2,
   breaklines=true,
   showtabs=false,
   captionpos=b
   % I listati sono oggetti mobili, cosi' non vengono mai spezzati fra due
   % pagine e la didascalia resta attaccata al codice. La chiave 'float' NON
   % ha effetto se data qui: va messa nell'argomento opzionale di ogni
   % lstlisting, e li' infatti c'e'. Un 'float=htbp' in questo blocco lascia
   % credere di coprire il caso senza coprirlo, quindi non c'e'.
}

% Nessuna riga vedova o orfana: una riga sola staccata dal suo paragrafo
% e' il difetto di impaginazione piu' visibile in un testo interlinea 1,5
\widowpenalty=10000
\clubpenalty=10000

% Nessun oggetto mobile scavalca l'inizio di una sottosezione. Senza questa
% barriera un listato o una figura che non entrano nella pagina corrente
% vengono rinviati alla prima posizione utile, che puo' cadere DOPO il titolo
% della sottosezione seguente: il codice di reduce finiva sotto il titolo di
% matmul, dove si legge come se fosse di matmul. L'opzione [section] copre i
% titoli di sezione, la ridefinizione sotto quelli di sottosezione.
\usepackage[section]{placeins}
\let\originalsubsection\subsection
\renewcommand{\subsection}{\FloatBarrier\originalsubsection}

% Didascalie: etichetta (Figura 3.1, Codice 3.2, Tabella 5.4) in grassetto e
% testo in corsivo, cosi' le didascalie lunghe si staccano dal corpo del testo.
% Caricato dopo 'listings' perche' valga anche per i lstlisting.
\usepackage{caption}
\captionsetup{labelfont=bf, textfont=it}

% Parametri dei float. Con i valori predefiniti una figura alta piu' di meta'
% pagina finisce su una pagina tutta sua, con fasce bianche sopra e sotto:
% questi permettono di mettere figura e testo sulla stessa pagina.
\renewcommand{\topfraction}{0.85}      % quota di pagina occupabile da float in alto
\renewcommand{\bottomfraction}{0.7}
\renewcommand{\textfraction}{0.1}      % testo minimo su una pagina mista
\renewcommand{\floatpagefraction}{0.8} % quanto deve essere pieno per meritare una pagina sua
\makeatletter
\setlength{\@fptop}{0pt}               % le pagine di soli float partono dall'alto
\makeatother

% Interruttore per le bozze parziali. Un file principale che scriva
% \conclusionifalse subito dopo % !TEX root = Tesi_Corsi_Leonardo_656276.tex
% Preambolo condiviso fra la tesi completa e il PDF dei soli capitoli 1-3.
% Non contiene \begin{document}: lo apre corpo.tex.

% Template Tesi in Informatica, ispirata al template "Tesi di laurea - Università di Pisa" di Simone Schirinzi (https://www.overleaf.com/latex/templates/tesi-di-laurea-universita-di-pisa/rwdcqtqwftpg).

% Carattere dimensione 12
\documentclass[12pt]{report}

% Per la stampa fronte-retro sostituire con:
% \documentclass[12pt, twoside]{report}

% Margini (4cm a sx, 2.5cm a dx, 2.5cm in alto, 2.5cm in basso)
% head=21.75pt: spazio per la testatina, il valore richiesto da scrlayer-scrpage
\usepackage[top=2.5cm, bottom=2.5cm, left=4cm, right=2.5cm, head=21.75pt, centering]{geometry}

% Per la stampa fronte-retro sostituire con: 
% \usepackage[top=2.5cm, bottom=2.5cm, inner=4cm, outer=4cm, right=2.5cm, centering]{geometry}

% Interlinea
\linespread{1.5}

% Librerie utili
\usepackage[italian]{babel} % applicazione regole di scrittura per la lingua italiana 
\usepackage[utf8]{inputenc} % codifica UTF-8
\usepackage{scrlayer-scrpage} % stili pagina: testatine e pie' di pagina
% Testatina: capitolo a sinistra, numero di pagina a destra, riga di separazione.
% Le pagine di apertura capitolo restano senza testatina (stile plain) e il
% numero di pagina lo tengono in fondo: e' la forma fra parentesi quadre.
\automark[section]{chapter}
\ihead[]{\headmark}
\chead[]{}
\ohead[]{\pagemark}
\ifoot[]{}
\cfoot[]{}
\ofoot[\pagemark]{}
\KOMAoptions{headsepline=true}
\pagestyle{scrheadings}
\usepackage{mathptmx} % font Times New Roman (simile)
\usepackage{graphicx} % inserimento di immagini
\usepackage{csquotes} % per le citazioni "in blocco"
\usepackage{url} % per \url{} nelle voci di bibliografia
%\usepackage[backend=biber, sorting=nty, ]{biblatex} % bibliografia con pacchetto biblatex (https://ctan.org/pkg/biblatex?lang=en)
\appto{\bibsetup}{\raggedright}

\usepackage{titlesec} % per la formattazione dei titoli delle sezioni, capitoli etc.
\usepackage{float} % per il posizionamento delle immagini
\usepackage{booktabs} % \toprule \midrule \bottomrule per le tabelle dei risultati
\usepackage{enumitem} % spaziatura degli elenchi: \begin{itemize}[topsep=..., itemsep=...]
\usepackage{needspace} % \needspace{N\baselineskip}: non spezzare un blocco a fine pagina
\usepackage{tikz} % per i disegni dell'architettura
\usetikzlibrary{arrows.meta, calc}
\usepackage{pgfplots} % per i grafici delle misure
\pgfplotsset{compat=1.18}
\usepackage{amsmath} % per le formule dei costi
\usepackage{amssymb} % \square a fine dimostrazione

\usepackage{listings} % per il codice di programmazione
% Fonte https://en.wikibooks.org/wiki/LaTeX/Source_Code_Listings. Per la lista di sintassi riconosciute.
\renewcommand{\lstlistingname}{Codice}% Listing -> Codice
\usepackage{xcolor}  % stile del codice
\definecolor{mygreen}{rgb}{0,0.6,0}
\definecolor{mygray}{rgb}{0.5,0.5,0.5}
\definecolor{mymauve}{rgb}{0.58,0,0.82}
\definecolor{darkgray}{rgb}{.4,.4,.4}
\definecolor{navy}{HTML}{000080}
\definecolor{purple}{rgb}{0.65, 0.12, 0.82}
\definecolor{codepurple}{rgb}{0.58,0,0.82}
\definecolor{backcolour}{rgb}{0.95,0.95,0.92}

% Stili configurabili del codice (lslisting) 
\lstset{ %
belowcaptionskip=0.5em,
backgroundcolor=\color{backcolour}, % choose the background color; you must add \usepackage{color} or \usepackage{xcolor}
basicstyle=\footnotesize, % the size of the fonts that are used for the code
breakatwhitespace=false, % sets if automatic breaks should only happen at whitespace
breaklines=true, % sets automatic line breaking
captionpos=b, % sets the caption-position to bottom
commentstyle=\color{mygreen}, % comment style
deletekeywords={...}, % if you want to delete keywords from the given language
escapeinside={\%*}{*)}, % if you want to add LaTeX within your code
extendedchars=true, % lets you use non-ASCII characters; for 8-bits encodings only, does not work with UTF-8
frame=single, % adds a frame around the code
keepspaces=true, % keeps spaces in text, useful for keeping indentation of code (possibly needs columns=flexible)
keywordstyle=\color{codepurple}, % keyword style
% language=Octave, % the language of the code
morekeywords={*,...}, % if you want to add more keywords to the set
numbers=left, % where to put the line-numbers; possible values are (none, left, right)
numbersep=5pt, % how far the line-numbers are from the code
numberstyle=\tiny\color{mygray}, % the style that is used for the line-numbers
rulecolor=\color{black}, % if not set, the frame-color may be changed on line-breaks within not-black text (e.g. comments (green here))
showspaces=false, % show spaces everywhere adding particular underscores; it overrides 'showstringspaces'
showstringspaces=false, % underline spaces within strings only
showtabs=false, % show tabs within strings adding particular underscores
stepnumber=1, % the step between two line-numbers. If it's 1, each line will be numbered
stringstyle=\color{mymauve}, % string literal style
tabsize=2, % sets default tabsize to 2 spaces
title=\lstname % show the filename of files included with \lstinputlisting; also try caption instead of title
}

% END of listing package 

% Esempio riconoscimento sintassi JavaScript 

\lstdefinelanguage{JavaScript}{
  keywords={typeof, new, true, false, catch, function, return, null, catch, switch, var, const, let, if, in, while, do, else, case, break},
  keywordstyle=\color{blue}\bfseries,
  ndkeywords={class, export, boolean, throw, implements, import, this},
  ndkeywordstyle=\color{darkgray}\bfseries,
  identifierstyle=\color{black},
  sensitive=false,
  comment=[l]{//},
  morecomment=[s]{/*}{*/},
  commentstyle=\color{mygreen}\ttfamily,
  stringstyle=\color{red}\ttfamily,
  morestring=[b]',
  morestring=[b]"
}

% Assembly RISC-V: le istruzioni rv32i usate dai kernel, piu' le quattro
% istruzioni di comunicazione dell'estensione, evidenziate a parte.
\lstdefinelanguage{RISCV}{
  keywords={add, addi, sub, and, andi, or, ori, xor, xori, sll, slli, srl, srli,
            sra, srai, slt, slti, sltu, sltiu, lui, auipc, lw, lh, lhu, lb, lbu,
            sw, sh, sb, beq, bne, blt, bge, bltu, bgeu, beqz, bnez, blez, bgez,
            j, jal, jalr, mv, li, la, ret, nop, ecall, mul},
  keywordstyle=\color{codepurple}\bfseries,
  ndkeywords={IN, OUT, ISRDY, SETRDY, NORD, EST, SUD, OVEST},
  ndkeywordstyle=\color{navy}\bfseries,
  identifierstyle=\color{black},
  sensitive=true,
  comment=[l]{\#},
  commentstyle=\color{mygreen}\ttfamily,
  morestring=[b]",
  stringstyle=\color{mymauve}\ttfamily
}

\lstset{
   extendedchars=true,
   basicstyle=\footnotesize\ttfamily,
   showstringspaces=false,
   showspaces=false,
   numbers=left,
   numberstyle=\footnotesize,
   numbersep=9pt,
   tabsize=2,
   breaklines=true,
   showtabs=false,
   captionpos=b
   % I listati sono oggetti mobili, cosi' non vengono mai spezzati fra due
   % pagine e la didascalia resta attaccata al codice. La chiave 'float' NON
   % ha effetto se data qui: va messa nell'argomento opzionale di ogni
   % lstlisting, e li' infatti c'e'. Un 'float=htbp' in questo blocco lascia
   % credere di coprire il caso senza coprirlo, quindi non c'e'.
}

% Nessuna riga vedova o orfana: una riga sola staccata dal suo paragrafo
% e' il difetto di impaginazione piu' visibile in un testo interlinea 1,5
\widowpenalty=10000
\clubpenalty=10000

% Nessun oggetto mobile scavalca l'inizio di una sottosezione. Senza questa
% barriera un listato o una figura che non entrano nella pagina corrente
% vengono rinviati alla prima posizione utile, che puo' cadere DOPO il titolo
% della sottosezione seguente: il codice di reduce finiva sotto il titolo di
% matmul, dove si legge come se fosse di matmul. L'opzione [section] copre i
% titoli di sezione, la ridefinizione sotto quelli di sottosezione.
\usepackage[section]{placeins}
\let\originalsubsection\subsection
\renewcommand{\subsection}{\FloatBarrier\originalsubsection}

% Didascalie: etichetta (Figura 3.1, Codice 3.2, Tabella 5.4) in grassetto e
% testo in corsivo, cosi' le didascalie lunghe si staccano dal corpo del testo.
% Caricato dopo 'listings' perche' valga anche per i lstlisting.
\usepackage{caption}
\captionsetup{labelfont=bf, textfont=it}

% Parametri dei float. Con i valori predefiniti una figura alta piu' di meta'
% pagina finisce su una pagina tutta sua, con fasce bianche sopra e sotto:
% questi permettono di mettere figura e testo sulla stessa pagina.
\renewcommand{\topfraction}{0.85}      % quota di pagina occupabile da float in alto
\renewcommand{\bottomfraction}{0.7}
\renewcommand{\textfraction}{0.1}      % testo minimo su una pagina mista
\renewcommand{\floatpagefraction}{0.8} % quanto deve essere pieno per meritare una pagina sua
\makeatletter
\setlength{\@fptop}{0pt}               % le pagine di soli float partono dall'alto
\makeatother

% Interruttore per le bozze parziali. Un file principale che scriva
% \conclusionifalse subito dopo % !TEX root = Tesi_Corsi_Leonardo_656276.tex
% Preambolo condiviso fra la tesi completa e il PDF dei soli capitoli 1-3.
% Non contiene \begin{document}: lo apre corpo.tex.

% Template Tesi in Informatica, ispirata al template "Tesi di laurea - Università di Pisa" di Simone Schirinzi (https://www.overleaf.com/latex/templates/tesi-di-laurea-universita-di-pisa/rwdcqtqwftpg).

% Carattere dimensione 12
\documentclass[12pt]{report}

% Per la stampa fronte-retro sostituire con:
% \documentclass[12pt, twoside]{report}

% Margini (4cm a sx, 2.5cm a dx, 2.5cm in alto, 2.5cm in basso)
% head=21.75pt: spazio per la testatina, il valore richiesto da scrlayer-scrpage
\usepackage[top=2.5cm, bottom=2.5cm, left=4cm, right=2.5cm, head=21.75pt, centering]{geometry}

% Per la stampa fronte-retro sostituire con: 
% \usepackage[top=2.5cm, bottom=2.5cm, inner=4cm, outer=4cm, right=2.5cm, centering]{geometry}

% Interlinea
\linespread{1.5}

% Librerie utili
\usepackage[italian]{babel} % applicazione regole di scrittura per la lingua italiana 
\usepackage[utf8]{inputenc} % codifica UTF-8
\usepackage{scrlayer-scrpage} % stili pagina: testatine e pie' di pagina
% Testatina: capitolo a sinistra, numero di pagina a destra, riga di separazione.
% Le pagine di apertura capitolo restano senza testatina (stile plain) e il
% numero di pagina lo tengono in fondo: e' la forma fra parentesi quadre.
\automark[section]{chapter}
\ihead[]{\headmark}
\chead[]{}
\ohead[]{\pagemark}
\ifoot[]{}
\cfoot[]{}
\ofoot[\pagemark]{}
\KOMAoptions{headsepline=true}
\pagestyle{scrheadings}
\usepackage{mathptmx} % font Times New Roman (simile)
\usepackage{graphicx} % inserimento di immagini
\usepackage{csquotes} % per le citazioni "in blocco"
\usepackage{url} % per \url{} nelle voci di bibliografia
%\usepackage[backend=biber, sorting=nty, ]{biblatex} % bibliografia con pacchetto biblatex (https://ctan.org/pkg/biblatex?lang=en)
\appto{\bibsetup}{\raggedright}

\usepackage{titlesec} % per la formattazione dei titoli delle sezioni, capitoli etc.
\usepackage{float} % per il posizionamento delle immagini
\usepackage{booktabs} % \toprule \midrule \bottomrule per le tabelle dei risultati
\usepackage{enumitem} % spaziatura degli elenchi: \begin{itemize}[topsep=..., itemsep=...]
\usepackage{needspace} % \needspace{N\baselineskip}: non spezzare un blocco a fine pagina
\usepackage{tikz} % per i disegni dell'architettura
\usetikzlibrary{arrows.meta, calc}
\usepackage{pgfplots} % per i grafici delle misure
\pgfplotsset{compat=1.18}
\usepackage{amsmath} % per le formule dei costi
\usepackage{amssymb} % \square a fine dimostrazione

\usepackage{listings} % per il codice di programmazione
% Fonte https://en.wikibooks.org/wiki/LaTeX/Source_Code_Listings. Per la lista di sintassi riconosciute.
\renewcommand{\lstlistingname}{Codice}% Listing -> Codice
\usepackage{xcolor}  % stile del codice
\definecolor{mygreen}{rgb}{0,0.6,0}
\definecolor{mygray}{rgb}{0.5,0.5,0.5}
\definecolor{mymauve}{rgb}{0.58,0,0.82}
\definecolor{darkgray}{rgb}{.4,.4,.4}
\definecolor{navy}{HTML}{000080}
\definecolor{purple}{rgb}{0.65, 0.12, 0.82}
\definecolor{codepurple}{rgb}{0.58,0,0.82}
\definecolor{backcolour}{rgb}{0.95,0.95,0.92}

% Stili configurabili del codice (lslisting) 
\lstset{ %
belowcaptionskip=0.5em,
backgroundcolor=\color{backcolour}, % choose the background color; you must add \usepackage{color} or \usepackage{xcolor}
basicstyle=\footnotesize, % the size of the fonts that are used for the code
breakatwhitespace=false, % sets if automatic breaks should only happen at whitespace
breaklines=true, % sets automatic line breaking
captionpos=b, % sets the caption-position to bottom
commentstyle=\color{mygreen}, % comment style
deletekeywords={...}, % if you want to delete keywords from the given language
escapeinside={\%*}{*)}, % if you want to add LaTeX within your code
extendedchars=true, % lets you use non-ASCII characters; for 8-bits encodings only, does not work with UTF-8
frame=single, % adds a frame around the code
keepspaces=true, % keeps spaces in text, useful for keeping indentation of code (possibly needs columns=flexible)
keywordstyle=\color{codepurple}, % keyword style
% language=Octave, % the language of the code
morekeywords={*,...}, % if you want to add more keywords to the set
numbers=left, % where to put the line-numbers; possible values are (none, left, right)
numbersep=5pt, % how far the line-numbers are from the code
numberstyle=\tiny\color{mygray}, % the style that is used for the line-numbers
rulecolor=\color{black}, % if not set, the frame-color may be changed on line-breaks within not-black text (e.g. comments (green here))
showspaces=false, % show spaces everywhere adding particular underscores; it overrides 'showstringspaces'
showstringspaces=false, % underline spaces within strings only
showtabs=false, % show tabs within strings adding particular underscores
stepnumber=1, % the step between two line-numbers. If it's 1, each line will be numbered
stringstyle=\color{mymauve}, % string literal style
tabsize=2, % sets default tabsize to 2 spaces
title=\lstname % show the filename of files included with \lstinputlisting; also try caption instead of title
}

% END of listing package 

% Esempio riconoscimento sintassi JavaScript 

\lstdefinelanguage{JavaScript}{
  keywords={typeof, new, true, false, catch, function, return, null, catch, switch, var, const, let, if, in, while, do, else, case, break},
  keywordstyle=\color{blue}\bfseries,
  ndkeywords={class, export, boolean, throw, implements, import, this},
  ndkeywordstyle=\color{darkgray}\bfseries,
  identifierstyle=\color{black},
  sensitive=false,
  comment=[l]{//},
  morecomment=[s]{/*}{*/},
  commentstyle=\color{mygreen}\ttfamily,
  stringstyle=\color{red}\ttfamily,
  morestring=[b]',
  morestring=[b]"
}

% Assembly RISC-V: le istruzioni rv32i usate dai kernel, piu' le quattro
% istruzioni di comunicazione dell'estensione, evidenziate a parte.
\lstdefinelanguage{RISCV}{
  keywords={add, addi, sub, and, andi, or, ori, xor, xori, sll, slli, srl, srli,
            sra, srai, slt, slti, sltu, sltiu, lui, auipc, lw, lh, lhu, lb, lbu,
            sw, sh, sb, beq, bne, blt, bge, bltu, bgeu, beqz, bnez, blez, bgez,
            j, jal, jalr, mv, li, la, ret, nop, ecall, mul},
  keywordstyle=\color{codepurple}\bfseries,
  ndkeywords={IN, OUT, ISRDY, SETRDY, NORD, EST, SUD, OVEST},
  ndkeywordstyle=\color{navy}\bfseries,
  identifierstyle=\color{black},
  sensitive=true,
  comment=[l]{\#},
  commentstyle=\color{mygreen}\ttfamily,
  morestring=[b]",
  stringstyle=\color{mymauve}\ttfamily
}

\lstset{
   extendedchars=true,
   basicstyle=\footnotesize\ttfamily,
   showstringspaces=false,
   showspaces=false,
   numbers=left,
   numberstyle=\footnotesize,
   numbersep=9pt,
   tabsize=2,
   breaklines=true,
   showtabs=false,
   captionpos=b
   % I listati sono oggetti mobili, cosi' non vengono mai spezzati fra due
   % pagine e la didascalia resta attaccata al codice. La chiave 'float' NON
   % ha effetto se data qui: va messa nell'argomento opzionale di ogni
   % lstlisting, e li' infatti c'e'. Un 'float=htbp' in questo blocco lascia
   % credere di coprire il caso senza coprirlo, quindi non c'e'.
}

% Nessuna riga vedova o orfana: una riga sola staccata dal suo paragrafo
% e' il difetto di impaginazione piu' visibile in un testo interlinea 1,5
\widowpenalty=10000
\clubpenalty=10000

% Nessun oggetto mobile scavalca l'inizio di una sottosezione. Senza questa
% barriera un listato o una figura che non entrano nella pagina corrente
% vengono rinviati alla prima posizione utile, che puo' cadere DOPO il titolo
% della sottosezione seguente: il codice di reduce finiva sotto il titolo di
% matmul, dove si legge come se fosse di matmul. L'opzione [section] copre i
% titoli di sezione, la ridefinizione sotto quelli di sottosezione.
\usepackage[section]{placeins}
\let\originalsubsection\subsection
\renewcommand{\subsection}{\FloatBarrier\originalsubsection}

% Didascalie: etichetta (Figura 3.1, Codice 3.2, Tabella 5.4) in grassetto e
% testo in corsivo, cosi' le didascalie lunghe si staccano dal corpo del testo.
% Caricato dopo 'listings' perche' valga anche per i lstlisting.
\usepackage{caption}
\captionsetup{labelfont=bf, textfont=it}

% Parametri dei float. Con i valori predefiniti una figura alta piu' di meta'
% pagina finisce su una pagina tutta sua, con fasce bianche sopra e sotto:
% questi permettono di mettere figura e testo sulla stessa pagina.
\renewcommand{\topfraction}{0.85}      % quota di pagina occupabile da float in alto
\renewcommand{\bottomfraction}{0.7}
\renewcommand{\textfraction}{0.1}      % testo minimo su una pagina mista
\renewcommand{\floatpagefraction}{0.8} % quanto deve essere pieno per meritare una pagina sua
\makeatletter
\setlength{\@fptop}{0pt}               % le pagine di soli float partono dall'alto
\makeatother

% Interruttore per le bozze parziali. Un file principale che scriva
% \conclusionifalse subito dopo % !TEX root = Tesi_Corsi_Leonardo_656276.tex
% Preambolo condiviso fra la tesi completa e il PDF dei soli capitoli 1-3.
% Non contiene \begin{document}: lo apre corpo.tex.

% Template Tesi in Informatica, ispirata al template "Tesi di laurea - Università di Pisa" di Simone Schirinzi (https://www.overleaf.com/latex/templates/tesi-di-laurea-universita-di-pisa/rwdcqtqwftpg).

% Carattere dimensione 12
\documentclass[12pt]{report}

% Per la stampa fronte-retro sostituire con:
% \documentclass[12pt, twoside]{report}

% Margini (4cm a sx, 2.5cm a dx, 2.5cm in alto, 2.5cm in basso)
% head=21.75pt: spazio per la testatina, il valore richiesto da scrlayer-scrpage
\usepackage[top=2.5cm, bottom=2.5cm, left=4cm, right=2.5cm, head=21.75pt, centering]{geometry}

% Per la stampa fronte-retro sostituire con: 
% \usepackage[top=2.5cm, bottom=2.5cm, inner=4cm, outer=4cm, right=2.5cm, centering]{geometry}

% Interlinea
\linespread{1.5}

% Librerie utili
\usepackage[italian]{babel} % applicazione regole di scrittura per la lingua italiana 
\usepackage[utf8]{inputenc} % codifica UTF-8
\usepackage{scrlayer-scrpage} % stili pagina: testatine e pie' di pagina
% Testatina: capitolo a sinistra, numero di pagina a destra, riga di separazione.
% Le pagine di apertura capitolo restano senza testatina (stile plain) e il
% numero di pagina lo tengono in fondo: e' la forma fra parentesi quadre.
\automark[section]{chapter}
\ihead[]{\headmark}
\chead[]{}
\ohead[]{\pagemark}
\ifoot[]{}
\cfoot[]{}
\ofoot[\pagemark]{}
\KOMAoptions{headsepline=true}
\pagestyle{scrheadings}
\usepackage{mathptmx} % font Times New Roman (simile)
\usepackage{graphicx} % inserimento di immagini
\usepackage{csquotes} % per le citazioni "in blocco"
\usepackage{url} % per \url{} nelle voci di bibliografia
%\usepackage[backend=biber, sorting=nty, ]{biblatex} % bibliografia con pacchetto biblatex (https://ctan.org/pkg/biblatex?lang=en)
\appto{\bibsetup}{\raggedright}

\usepackage{titlesec} % per la formattazione dei titoli delle sezioni, capitoli etc.
\usepackage{float} % per il posizionamento delle immagini
\usepackage{booktabs} % \toprule \midrule \bottomrule per le tabelle dei risultati
\usepackage{enumitem} % spaziatura degli elenchi: \begin{itemize}[topsep=..., itemsep=...]
\usepackage{needspace} % \needspace{N\baselineskip}: non spezzare un blocco a fine pagina
\usepackage{tikz} % per i disegni dell'architettura
\usetikzlibrary{arrows.meta, calc}
\usepackage{pgfplots} % per i grafici delle misure
\pgfplotsset{compat=1.18}
\usepackage{amsmath} % per le formule dei costi
\usepackage{amssymb} % \square a fine dimostrazione

\usepackage{listings} % per il codice di programmazione
% Fonte https://en.wikibooks.org/wiki/LaTeX/Source_Code_Listings. Per la lista di sintassi riconosciute.
\renewcommand{\lstlistingname}{Codice}% Listing -> Codice
\usepackage{xcolor}  % stile del codice
\definecolor{mygreen}{rgb}{0,0.6,0}
\definecolor{mygray}{rgb}{0.5,0.5,0.5}
\definecolor{mymauve}{rgb}{0.58,0,0.82}
\definecolor{darkgray}{rgb}{.4,.4,.4}
\definecolor{navy}{HTML}{000080}
\definecolor{purple}{rgb}{0.65, 0.12, 0.82}
\definecolor{codepurple}{rgb}{0.58,0,0.82}
\definecolor{backcolour}{rgb}{0.95,0.95,0.92}

% Stili configurabili del codice (lslisting) 
\lstset{ %
belowcaptionskip=0.5em,
backgroundcolor=\color{backcolour}, % choose the background color; you must add \usepackage{color} or \usepackage{xcolor}
basicstyle=\footnotesize, % the size of the fonts that are used for the code
breakatwhitespace=false, % sets if automatic breaks should only happen at whitespace
breaklines=true, % sets automatic line breaking
captionpos=b, % sets the caption-position to bottom
commentstyle=\color{mygreen}, % comment style
deletekeywords={...}, % if you want to delete keywords from the given language
escapeinside={\%*}{*)}, % if you want to add LaTeX within your code
extendedchars=true, % lets you use non-ASCII characters; for 8-bits encodings only, does not work with UTF-8
frame=single, % adds a frame around the code
keepspaces=true, % keeps spaces in text, useful for keeping indentation of code (possibly needs columns=flexible)
keywordstyle=\color{codepurple}, % keyword style
% language=Octave, % the language of the code
morekeywords={*,...}, % if you want to add more keywords to the set
numbers=left, % where to put the line-numbers; possible values are (none, left, right)
numbersep=5pt, % how far the line-numbers are from the code
numberstyle=\tiny\color{mygray}, % the style that is used for the line-numbers
rulecolor=\color{black}, % if not set, the frame-color may be changed on line-breaks within not-black text (e.g. comments (green here))
showspaces=false, % show spaces everywhere adding particular underscores; it overrides 'showstringspaces'
showstringspaces=false, % underline spaces within strings only
showtabs=false, % show tabs within strings adding particular underscores
stepnumber=1, % the step between two line-numbers. If it's 1, each line will be numbered
stringstyle=\color{mymauve}, % string literal style
tabsize=2, % sets default tabsize to 2 spaces
title=\lstname % show the filename of files included with \lstinputlisting; also try caption instead of title
}

% END of listing package 

% Esempio riconoscimento sintassi JavaScript 

\lstdefinelanguage{JavaScript}{
  keywords={typeof, new, true, false, catch, function, return, null, catch, switch, var, const, let, if, in, while, do, else, case, break},
  keywordstyle=\color{blue}\bfseries,
  ndkeywords={class, export, boolean, throw, implements, import, this},
  ndkeywordstyle=\color{darkgray}\bfseries,
  identifierstyle=\color{black},
  sensitive=false,
  comment=[l]{//},
  morecomment=[s]{/*}{*/},
  commentstyle=\color{mygreen}\ttfamily,
  stringstyle=\color{red}\ttfamily,
  morestring=[b]',
  morestring=[b]"
}

% Assembly RISC-V: le istruzioni rv32i usate dai kernel, piu' le quattro
% istruzioni di comunicazione dell'estensione, evidenziate a parte.
\lstdefinelanguage{RISCV}{
  keywords={add, addi, sub, and, andi, or, ori, xor, xori, sll, slli, srl, srli,
            sra, srai, slt, slti, sltu, sltiu, lui, auipc, lw, lh, lhu, lb, lbu,
            sw, sh, sb, beq, bne, blt, bge, bltu, bgeu, beqz, bnez, blez, bgez,
            j, jal, jalr, mv, li, la, ret, nop, ecall, mul},
  keywordstyle=\color{codepurple}\bfseries,
  ndkeywords={IN, OUT, ISRDY, SETRDY, NORD, EST, SUD, OVEST},
  ndkeywordstyle=\color{navy}\bfseries,
  identifierstyle=\color{black},
  sensitive=true,
  comment=[l]{\#},
  commentstyle=\color{mygreen}\ttfamily,
  morestring=[b]",
  stringstyle=\color{mymauve}\ttfamily
}

\lstset{
   extendedchars=true,
   basicstyle=\footnotesize\ttfamily,
   showstringspaces=false,
   showspaces=false,
   numbers=left,
   numberstyle=\footnotesize,
   numbersep=9pt,
   tabsize=2,
   breaklines=true,
   showtabs=false,
   captionpos=b
   % I listati sono oggetti mobili, cosi' non vengono mai spezzati fra due
   % pagine e la didascalia resta attaccata al codice. La chiave 'float' NON
   % ha effetto se data qui: va messa nell'argomento opzionale di ogni
   % lstlisting, e li' infatti c'e'. Un 'float=htbp' in questo blocco lascia
   % credere di coprire il caso senza coprirlo, quindi non c'e'.
}

% Nessuna riga vedova o orfana: una riga sola staccata dal suo paragrafo
% e' il difetto di impaginazione piu' visibile in un testo interlinea 1,5
\widowpenalty=10000
\clubpenalty=10000

% Nessun oggetto mobile scavalca l'inizio di una sottosezione. Senza questa
% barriera un listato o una figura che non entrano nella pagina corrente
% vengono rinviati alla prima posizione utile, che puo' cadere DOPO il titolo
% della sottosezione seguente: il codice di reduce finiva sotto il titolo di
% matmul, dove si legge come se fosse di matmul. L'opzione [section] copre i
% titoli di sezione, la ridefinizione sotto quelli di sottosezione.
\usepackage[section]{placeins}
\let\originalsubsection\subsection
\renewcommand{\subsection}{\FloatBarrier\originalsubsection}

% Didascalie: etichetta (Figura 3.1, Codice 3.2, Tabella 5.4) in grassetto e
% testo in corsivo, cosi' le didascalie lunghe si staccano dal corpo del testo.
% Caricato dopo 'listings' perche' valga anche per i lstlisting.
\usepackage{caption}
\captionsetup{labelfont=bf, textfont=it}

% Parametri dei float. Con i valori predefiniti una figura alta piu' di meta'
% pagina finisce su una pagina tutta sua, con fasce bianche sopra e sotto:
% questi permettono di mettere figura e testo sulla stessa pagina.
\renewcommand{\topfraction}{0.85}      % quota di pagina occupabile da float in alto
\renewcommand{\bottomfraction}{0.7}
\renewcommand{\textfraction}{0.1}      % testo minimo su una pagina mista
\renewcommand{\floatpagefraction}{0.8} % quanto deve essere pieno per meritare una pagina sua
\makeatletter
\setlength{\@fptop}{0pt}               % le pagine di soli float partono dall'alto
\makeatother

% Interruttore per le bozze parziali. Un file principale che scriva
% \conclusionifalse subito dopo \input{preambolo} esclude il capitolo 6 dalla
% composizione, dall'indice e dai rimandi che lo nominano, in un colpo solo.
% Va fatto cosi' e non con \includeonly: per un capitolo escluso da
% \includeonly LaTeX legge lo stesso il suo .aux, quindi le voci finirebbero
% comunque nell'indice e i contatori di fine capitolo farebbero saltare la
% numerazione delle pagine.
\newif\ifconclusioni
\conclusionitrue

% Formato delle intestazioni
% Intestazione di capitolo su due righe: "Capitolo N" e sotto il titolo.
% \chaptertitlename vale "Capitolo" e diventa "Appendice" dopo \appendix.
\titleformat{\chapter}[display]
  {\normalfont\huge\bfseries}   % formato del blocco
  {\chaptertitlename\ \thechapter} % etichetta
  {20pt}                        % distanza fra etichetta e titolo
  {\Huge}                       % formato del titolo
\titlespacing*{\chapter}{0pt}{50pt}{40pt}

% Indice fino alle sottosezioni: i \paragraph restano fuori
\setcounter{tocdepth}{2}

% Sorgenti delle appendici: listings non gestisce da solo le lettere accentate
% in UTF-8, e conosce Verilog ma non le parole chiave di SystemVerilog
\lstset{literate={à}{{\`a}}1 {è}{{\`e}}1 {é}{{\'e}}1 {È}{{\`E}}1
  {ì}{{\`i}}1 {ò}{{\`o}}1 {ù}{{\`u}}1 {±}{{$\pm$}}1 {—}{{\textemdash}}1}
\lstdefinelanguage[System]{Verilog}[]{Verilog}{morekeywords={logic, always_ff,
  always_comb, localparam, genvar, generate, endgenerate}}

% Le regole di sillabazione italiane spezzavano il nome in Sy-stemVerilog
\babelhyphenation[italian]{System-Verilog}
 esclude il capitolo 6 dalla
% composizione, dall'indice e dai rimandi che lo nominano, in un colpo solo.
% Va fatto cosi' e non con \includeonly: per un capitolo escluso da
% \includeonly LaTeX legge lo stesso il suo .aux, quindi le voci finirebbero
% comunque nell'indice e i contatori di fine capitolo farebbero saltare la
% numerazione delle pagine.
\newif\ifconclusioni
\conclusionitrue

% Formato delle intestazioni
% Intestazione di capitolo su due righe: "Capitolo N" e sotto il titolo.
% \chaptertitlename vale "Capitolo" e diventa "Appendice" dopo \appendix.
\titleformat{\chapter}[display]
  {\normalfont\huge\bfseries}   % formato del blocco
  {\chaptertitlename\ \thechapter} % etichetta
  {20pt}                        % distanza fra etichetta e titolo
  {\Huge}                       % formato del titolo
\titlespacing*{\chapter}{0pt}{50pt}{40pt}

% Indice fino alle sottosezioni: i \paragraph restano fuori
\setcounter{tocdepth}{2}

% Sorgenti delle appendici: listings non gestisce da solo le lettere accentate
% in UTF-8, e conosce Verilog ma non le parole chiave di SystemVerilog
\lstset{literate={à}{{\`a}}1 {è}{{\`e}}1 {é}{{\'e}}1 {È}{{\`E}}1
  {ì}{{\`i}}1 {ò}{{\`o}}1 {ù}{{\`u}}1 {±}{{$\pm$}}1 {—}{{\textemdash}}1}
\lstdefinelanguage[System]{Verilog}[]{Verilog}{morekeywords={logic, always_ff,
  always_comb, localparam, genvar, generate, endgenerate}}

% Le regole di sillabazione italiane spezzavano il nome in Sy-stemVerilog
\babelhyphenation[italian]{System-Verilog}
 esclude il capitolo 6 dalla
% composizione, dall'indice e dai rimandi che lo nominano, in un colpo solo.
% Va fatto cosi' e non con \includeonly: per un capitolo escluso da
% \includeonly LaTeX legge lo stesso il suo .aux, quindi le voci finirebbero
% comunque nell'indice e i contatori di fine capitolo farebbero saltare la
% numerazione delle pagine.
\newif\ifconclusioni
\conclusionitrue

% Formato delle intestazioni
% Intestazione di capitolo su due righe: "Capitolo N" e sotto il titolo.
% \chaptertitlename vale "Capitolo" e diventa "Appendice" dopo \appendix.
\titleformat{\chapter}[display]
  {\normalfont\huge\bfseries}   % formato del blocco
  {\chaptertitlename\ \thechapter} % etichetta
  {20pt}                        % distanza fra etichetta e titolo
  {\Huge}                       % formato del titolo
\titlespacing*{\chapter}{0pt}{50pt}{40pt}

% Indice fino alle sottosezioni: i \paragraph restano fuori
\setcounter{tocdepth}{2}

% Sorgenti delle appendici: listings non gestisce da solo le lettere accentate
% in UTF-8, e conosce Verilog ma non le parole chiave di SystemVerilog
\lstset{literate={à}{{\`a}}1 {è}{{\`e}}1 {é}{{\'e}}1 {È}{{\`E}}1
  {ì}{{\`i}}1 {ò}{{\`o}}1 {ù}{{\`u}}1 {±}{{$\pm$}}1 {—}{{\textemdash}}1}
\lstdefinelanguage[System]{Verilog}[]{Verilog}{morekeywords={logic, always_ff,
  always_comb, localparam, genvar, generate, endgenerate}}

% Le regole di sillabazione italiane spezzavano il nome in Sy-stemVerilog
\babelhyphenation[italian]{System-Verilog}
 esclude il capitolo 6 dalla
% composizione, dall'indice e dai rimandi che lo nominano, in un colpo solo.
% Va fatto cosi' e non con \includeonly: per un capitolo escluso da
% \includeonly LaTeX legge lo stesso il suo .aux, quindi le voci finirebbero
% comunque nell'indice e i contatori di fine capitolo farebbero saltare la
% numerazione delle pagine.
\newif\ifconclusioni
\conclusionitrue

% Formato delle intestazioni
% Intestazione di capitolo su due righe: "Capitolo N" e sotto il titolo.
% \chaptertitlename vale "Capitolo" e diventa "Appendice" dopo \appendix.
\titleformat{\chapter}[display]
  {\normalfont\huge\bfseries}   % formato del blocco
  {\chaptertitlename\ \thechapter} % etichetta
  {20pt}                        % distanza fra etichetta e titolo
  {\Huge}                       % formato del titolo
\titlespacing*{\chapter}{0pt}{50pt}{40pt}

% Indice fino alle sottosezioni: i \paragraph restano fuori
\setcounter{tocdepth}{2}

% Sorgenti delle appendici: listings non gestisce da solo le lettere accentate
% in UTF-8, e conosce Verilog ma non le parole chiave di SystemVerilog
\lstset{literate={à}{{\`a}}1 {è}{{\`e}}1 {é}{{\'e}}1 {È}{{\`E}}1
  {ì}{{\`i}}1 {ò}{{\`o}}1 {ù}{{\`u}}1 {±}{{$\pm$}}1 {—}{{\textemdash}}1}
\lstdefinelanguage[System]{Verilog}[]{Verilog}{morekeywords={logic, always_ff,
  always_comb, localparam, genvar, generate, endgenerate}}

% Le regole di sillabazione italiane spezzavano il nome in Sy-stemVerilog
\babelhyphenation[italian]{System-Verilog}
